%iffalse
\let\negmedspace\undefined
\let\negthickspace\undefined
\documentclass[journal,12pt,onecolumn]{IEEEtran}
\usepackage{cite}
\usepackage{amsmath,amssymb,amsfonts,amsthm}
\usepackage{algorithmic}
\usepackage{graphicx}
\usepackage{textcomp}
\usepackage{xcolor}
\usepackage{txfonts}
\usepackage{listings}
\usepackage{enumitem}
\usepackage{mathtools}
\usepackage{gensymb}
\usepackage{comment}
\usepackage[breaklinks=true]{hyperref}
\usepackage{tkz-euclide}
\usepackage{listings}
\usepackage{gvv}
%\def\inputGnumericTable{}
\usepackage[latin1]{inputenc}
\usepackage{color}
\usepackage{array}
\usepackage{longtable}
\usepackage{calc}
\usepackage{multirow}
\usepackage{hhline}
\usepackage{ifthen}
\usepackage{lscape}
\usepackage{tabularx}
\usepackage{array}
\usepackage{float}
\usepackage{caption}
\usepackage{multicol}

\newtheorem{theorem}{Theorem}[section]
\newtheorem{problem}{Problem}
\newtheorem{proposition}{Proposition}[section]
\newtheorem{lemma}{Lemma}[section]
\newtheorem{corollary}[theorem]{Corollary}
\newtheorem{example}{Example}[section]
\newtheorem{definition}[problem]{Definition}
\newcommand{\BEQA}{\begin{eqnarray}}
\newcommand{\EEQA}{\end{eqnarray}}
%\newcommand{\define}{\stackrel{\triangle}{=}}
\theoremstyle{remark}
%\newtheorem{rem}{Remark}

% Marks the beginning of the document
\begin{document}
\bibliographystyle{IEEEtran}
\vspace{3cm}

\title{2.5.24}
\author{EE25BTECH11058 - Mohana Krishna Sushma}
\maketitle

\renewcommand{\thefigure}{\theenumi}
\renewcommand{\thetable}{\theenumi}

\textbf{Question:}\par
If $|\vec{a}| = 2$, $|\vec{b}| = 7$ and $\vec{a} \times \vec{b} = 3\hat{\imath}
+ 2\hat{\jmath} + 6\hat{k}$, find the angle between $\vec{a}$ and $\vec{b}$.
\par

\textbf{Solution:}\par

Let the vector $\vec{c} = \vec{a} \times \vec{b}$. We are given:
\begin{align}
    \norm{\vec{a}} &= 2 \\
    \norm{\vec{b}} &= 7 \\
    \vec{c} &= \myvec{3 \\ 2 \\ 6}
\end{align}

First, we calculate the squared magnitude of $\vec{c}$ using the matrix product $\vec{c}^\top\vec{c}$:
\begin{align}
    \norm{\vec{c}}^2 &= \vec{c}^\top\vec{c} \\
    &= \myvec{3 & 2 & 6} \myvec{3 \\ 2 \\ 6} \\
    &= (3)(3) + (2)(2) + (6)(6) \\
    &= 9 + 4 + 36 \\
    &= 49
\end{align}
\hspace{1cm}Therefore, the magnitude of $\vec{c}$ is $\norm{\vec{c}} = \sqrt{49} = 7$.

Using Lagrange's identity , 
\begin{align}
    \norm{\vec{a} \times \vec{b}}^2 + (\vec{a}^\top\vec{b})^2 = \norm{\vec{a}}^2\norm{\vec{b}}^2
\end{align}

Substituting the magnitudes we have:
\begin{align}
    \norm{\vec{c}}^2 + (\vec{a}^\top\vec{b})^2 &= (2)^2(7)^2 \\
    49 + (\vec{a}^\top\vec{b})^2 &= (4)(49) \\
    49 + (\vec{a}^\top\vec{b})^2 &= 196
\end{align}

Solving for the squared dot product:
\begin{align}
    (\vec{a}^\top\vec{b})^2 = 196 - 49 = 147
\end{align}

\begin{align}
    \vec{a}^\top\vec{b} = \sqrt{147} = \sqrt{49 \times 3} = 7\sqrt{3}
\end{align}

The definition of the dot product in terms of the angle $\theta$ between the vectors:
\begin{align}
    \vec{a}^\top\vec{b} = \norm{\vec{a}}\norm{\vec{b}}\cos(\theta)
\end{align}
Substituting the known values:
\begin{align}
    7\sqrt{3} &= (2)(7)\cos(\theta) \\
    7\sqrt{3} &= 14\cos(\theta)
\end{align}

Solving for $\cos(\theta)$:
\begin{align}
    \cos(\theta) = \frac{7\sqrt{3}}{14} = \frac{\sqrt{3}}{2}
\end{align}

This implies that the angle $\theta$ is:
\begin{align}
    \theta = \frac{\pi}{6} \text{ or } 30\degree
\end{align}

\begin{figure}[H]
    \centering
   \includegraphics[width=1\linewidth]{figs/fig.png}
    \label{fig:placeholder}
\end{figure}

\end{document}